\documentclass[10pt,a4paper]{amsart}
\usepackage[margin=19mm]{geometry}
\usepackage{amsmath,amssymb,mathtools}
\usepackage[scr=rsfs,cal=euler]{mathalfa}
\usepackage{xfrac}
\usepackage[unicode,hidelinks,bookmarks=false]{hyperref}
\newcommand{\ie}{i.e.\ }
\newcommand{\cf}{cf.\ }
\newcommand{\resp}{respectively\ }
\newcommand{\Subsection}{\subsection}
\providecommand{\nobreakdash}{\nobreak\hbox{-}}
\setlength{\emergencystretch}{2em}
\multlinegap0pt
\usepackage{amssymb}
\newtheorem{theorem}{Theorem}
\newtheorem{proposition}{Proposition}
\newtheorem{pf}{Pf}
\title{On the behavior of analytic torsion for twisted canonical bundles under degenerations}
\author{Ken-Ichi Yoshikawa}
\date{Source: arXiv:2602.14797v2 (CC BY 4.0)}
\begin{document}
\maketitle

\noindent\emph{Source and licence.} On the behavior of analytic torsion for twisted canonical bundles under degenerations. Version arXiv:2602.14797v2 (CC BY 4.0), distributed by arXiv under the Creative Commons Attribution 4.0 International licence, http://creativecommons.org/licenses/by/4.0/, as recorded in the arXiv licence metadata for this exact version and shown on the abstract page https://arxiv.org/abs/2602.14797v2. Full licence notice: \texttt{licenses/arxiv-2602.14797-CC-BY-4.0.txt}.\par\smallskip

\noindent\textit{Editorial scope.} Counted unit: the article's theorem thm:leading:term:torsion (source0.tex lines 3994--4004). The article's complete original proof of it is delivered with it (source0.tex lines 4006--4063, one \texttt{proof} environment of 58 lines). No mathematics is written, repaired or re-derived: the statement and the proof are the article's own lines, in the article's own notation, and no retained line is substituted or re-flowed.  Retained uncounted as the source-local context the proof needs: lines 2203--2257; lines 2483--2486; lines 3054--3064; lines 3109--3118; lines 3967--3975.  Nothing was omitted from any retained span, so the package carries no omission record.  No retained line contains a citation, so the wrapper carries no bibliography and no bibliography entry.  No retained line contains a figure environment, an image-inclusion command or drawing code, so no image is delivered and the package declares no asset dependency.  Editorial formatting only, in addition: an AMS \texttt{amsart} preamble that restates the article's own declarations and macros used by the retained lines, namely the environments \texttt{theorem}, \texttt{proposition}, \texttt{pf} on separate counters, together with the standard packages those lines need.  The retained environments print in this document as Theorem 1, Proposition 1, Pf 1 in order of appearance, whereas the article prints the same items with the numbering of its own full document.  The complete native source of the article as distributed in the arXiv e-print for this exact version is bundled unchanged as \texttt{sources/194744/source0.tex}.
\begin{theorem}
\label{thm:str:sing:L2}
By choosing suitable bases $\{ {\mathbf e}_{1},\ldots,{\mathbf e}_{h_{W}^{q}} \}$ of $R^{q}\pi_{*}K_{X/S}(\xi)_{W}$ and 
$\{ \widetilde{\mathbf e}_{1}, \ldots, \widetilde{\mathbf e}_{h_{W}^{q}}\}$ of $R^{q}f_{*}K_{Y/T}(F^{*}\xi)_{W}$,
there exist integers $e^{q}_{1},\ldots, e^{q}_{h_{W}^{q}} \geq0$ with the following properties:
\par{\rm (1) }
The $(h_{W}^{q}, h_{W}^{q})$-Hermitian matrices $H(s) := ( H_{\alpha\beta}(s) )$, 
$H_{\alpha\beta}(s) := ( {\mathbf e}_{\alpha} , {\mathbf e}_{\beta} )_{L^{2},X_{s}}$, and
$\widetilde{H}(t) = ( \widetilde{H}_{\alpha\beta}(t) )$, 
$\widetilde{H}_{\alpha\beta}(t) := ( \widetilde{\mathbf e}_{\alpha} , \widetilde{\mathbf e}_{\beta} )_{L^{2},Y_{t}}$, 
are related via the following identity
\begin{equation}
\label{eqn:decomp:H}
H(\mu(t)) = D(t)\cdot \widetilde{H}(t)\cdot\overline{D(t)},
\qquad
D(t) = {\rm diag}( t^{-e^{q}_{1}},\ldots,t^{-e^{q}_{h_{W}^{q}}} ).
\end{equation}
Moreover, $\widetilde{H}(t)$ admits an asymptotic expansion of Barlet-Takayama type
\begin{equation}
\label{eqn:asym:exp:H:2}
\widetilde{H}(t) \equiv \sum_{0\leq m \leq n} (\log |t|^{-2})^{m} A_{m} 
\mod \bigoplus_{0\leq k \leq n} (\log |t|^{-2})^{k} {\mathfrak m}^{\infty}_{0}(T) \otimes M_{h_{W}^{q}}( {\mathbf C} )
\end{equation}
with constant $(h_{W}^{q}, h_{W}^{q})$-Hermitian matrices $A_{m}$ $(1\leq m \leq n)$. 
In particular, there exist constants $a_{m}^{q} \in {\mathbf R}$ $( 1 \leq m \leq n h_{W}^{q})$ such that 
\begin{equation}
\label{eqn:asym:exp:det:H:2}
\det \widetilde{H}(t) = \sum_{0 \leq m \leq n h_{W}^{q}} a_{m}^{q} (\log |t|^{-2})^{m} 
\mod \bigoplus_{0\leq k \leq n h_{W}^{q}} (\log |t|^{-2})^{k} {\mathfrak m}^{\infty}_{0}(T).
\end{equation}
\par{\rm (2) }
There exists a constant $C>0$ such that $\widetilde{H}(t) \geq C\,I_{h_{W}^{q}}$ for all $t\in T^{o}$ as positive definite Hermitian matrices.
Especially, $a_{m}^{q} \not=0$ for some $1 \leq m \leq n h_{W}^{q}$.
\par{\rm (3) }
There exist real-valued smooth functions $\psi_{j,k}^{q}(s)$ on $S$ such that
\begin{equation}
\label{eqn:L2:asym:eqv}
\| {\mathbf e}_{1}\wedge \cdots\wedge {\mathbf e}_{h_{W}^{q}} \|_{L^{2}}^{2} =
|s|^{-2\delta_{W}^{q}} \sum_{0\leq m \leq n h_{W}^{q}} \{ a_{m}^{q} 
+ \sum_{1 \leq j \leq d} |s|^{\frac{2j}{d}} \psi_{j,m}^{q}(s) \} (\log |s|^{-2})^{m}. 
\end{equation}
Let
\begin{equation}
\label{eqn:rho:q}
\varrho_{W}^{q} := \max\{ 0 \leq m \leq n h_{W}^{q};\, a_{m}^{q} \not=0 \}.
\end{equation}
Then the following equality holds as $s\to 0$:
\begin{equation}
\label{eqn:asymptotoc:L2:X}
\log\| {\mathbf e}_{1}\wedge \cdots\wedge {\mathbf e}_{h_{W}^{q}} \|_{L^{2}}^{2} 
= -\delta_{W}^{q}\,\log |s|^{2} + \varrho_{W}^{q}\log\log(|s|^{-2}) + c_{W}^{q} + O\left(1/\log |s|^{-1} \right).
\end{equation}
\par{\rm (4) }
The rational numbers $e^{q}_{\alpha}/\deg \mu$ ($1 \leq \alpha \leq h_{W}^{q}$) are independent of the choice of semi-stable reduction.
\end{theorem}

\begin{equation}
\label{eqn:asymp:eq:tors}
\log \tau_{G}(X_{s},K_{X_{s}}(\xi_{s}))(g) = \kappa_{g} \log |s|^{2} - \varrho_{g}\,\log\log( |s|^{-2}) + \gamma_{g} + O\left(1/\log |s|^{-1} \right).
\end{equation}

\begin{proposition}
\label{prop:L2:log:can}
Suppose that $X_{0}$ is reduced and the pair $(X,X_{0})$ is log-canonical. 
Then the elementary exponents of $R^{q}\pi_{*}K_{X/S}(\xi)$ vanish.
Moreover, there exist constants $\varrho^{q} \in {\mathbf Z}_{\geq 0}$ and $c^{q} \in {\mathbf R}$ such that as $s\to0$,
\begin{equation}
\label{eqn:asymptotoc:L2:X:lc}
\log\det ( g_{\alpha\bar{\beta}}(s) )
= \varrho^{q} \log\log(|s|^{-2}) + c^{q} + O(1/\log |s|^{-1}).
\end{equation}
\end{proposition}

\begin{theorem}
\label{thm:sing:tor:log:can}
Suppose that $X_{0}$ is reduced and the pair $(X, X_{0})$ has only log-canonical singularities. 
Then there exist $\gamma_{g} \in {\mathbf C}$ such that as $s \to 0$,
$$
\log\tau_{G}(X_{s},K_{X_{s}}(\xi_{s}))(g) 
= \alpha_{\pi,K_{X/S}(\xi)}(g) \log |s|^{2} - \varrho_{g}\,\log\log( |s|^{-2}) + \gamma_{g} + O\left(1/\log |s|^{-1} \right).
$$
In particular, if $X_{0}$ is a reduced normal crossing divisor of $X$, this formula holds.
\end{theorem}

\begin{proposition}
\label{prop:comparison:quillen}
The following equality of functions on $T^{o}$ holds:
$$
\log \left( \frac{\|\cdot\|_{Q,\lambda( K_{Y/T}(F^{*}\xi) )}}{\|\cdot\|'_{Q,\lambda( K_{Y/T}(F^{*}\xi) )}} \right)^{2} 
\equiv_{\mathcal B} -\beta_{Y/X, \xi} \log|t|^{2}.
$$
In particular, $\beta_{Y/X, \xi}$ is independent of the choices of $M$ and $M^{\natural}$.
\end{proposition}

\begin{theorem}
\label{thm:leading:term:torsion}
The following equality of rational numbers holds:
\begin{equation}
\label{eqn:kappa}
\kappa_{\pi, K_{X/S}(\xi)} = \frac{1}{\deg\mu} \left( \alpha_{f, K_{Y/T}(F^{*}\xi)} + \beta_{Y/X, \xi } \right).
\end{equation}
In particular, $\kappa_{\pi, K_{X/S}(\xi)}$ is given by the integral of various characteristic classes associated to 
the semi-stable reduction of $\pi \colon X \to S$.
Moreover, the right hand side of \eqref{eqn:kappa} is independent of the choice of a semi-stable reduction of $\pi \colon X \to S$.
\end{theorem}

\begin{pf}
Let $\|\cdot\|_{L^{2},\lambda(K_{Y/T}(F^{*}\xi) )}$ be the $L^{2}$-metric on $\lambda(K_{Y/T}(F^{*}\xi))$ with respect to the metrics 
$h_{{\mathcal Y}/C'}$ and $F^{*}h_{\xi}$. Similarly, let $\|\cdot\|'_{L^{2},\lambda(K_{Y/T}(F^{*}\xi))}$ be the $L^{2}$-metric on 
$\lambda(K_{Y/T}(F^{*}\xi))$ with respect to the metrics $F^{*}h_{{\mathcal X}/C}$ and $F^{*}h_{\xi}$.
Let $\tau(Y_{t}, K_{Y_{t}}(F^{*}\xi))$ be the analytic torsion of $F^{*}\xi_{t}$ with respect to the metrics $h_{\mathcal Y}|_{Y_{t}}$ and 
$F^{*}h_{\xi}|_{Y_{t}}$.
On the other hand, the analytic torsion of $F^{*}\xi_{t}$ with respect to the metrics $F^{*}h_{{\mathcal X}/C}|_{Y_{t}}$ and 
$F^{*}h_{\xi}|_{Y_{t}}$ is given by $\tau(X_{\mu(t)}, K_{X_{\mu(t)}}(\xi))$. Hence 
\begin{equation}
\label{eqn:comparison:Quillen:1}
\log\left(\frac{\|\cdot\|_{Q,\lambda(K_{Y_{t}}(F^{*}\xi))}}{\|\cdot\|'_{Q,\lambda(K_{Y_{t}}(F^{*}\xi))}}\right)^{2}
=
\log\frac{\tau(Y_{t}, K_{Y_{t}}(F^{*}\xi))}{\tau(X_{\mu(t)}, K_{X_{\mu(t)}}(\xi))}
+
\log\left(\frac{\|\cdot\|_{L^{2},\lambda(K_{Y_{t}}(F^{*}\xi))}}{\|\cdot\|'_{L^{2},\lambda(K_{Y_{t}}(F^{*}\xi))}}\right)^{2}.
\end{equation}
Since $Y_{0}$ is reduced and normal crossing, it follows from Theorem~\ref{thm:sing:tor:log:can} that as $t \to 0$,
\begin{equation}
\label{eqn:asymptotics:torsion:1}
\log\tau(Y_{t}, K_{Y_{t}}(F^{*}\xi_{t})) = \alpha_{f, K_{Y/T}(F^{*}\xi)} \log|t|^{2} + O\left( \log\log (| t |^{-2}) \right).
\end{equation}
\par
Let $\sigma$ be a nowhere vanishing holomorphic section of $\lambda(K_{Y/T}(F^{*}\xi))$ defined on $T$. 
By Proposition~\ref{prop:L2:log:can}, as $t\to0$, we have
\begin{equation}
\label{eqn:estimate:L2:metric:1}
\log\|\sigma(t)\|_{L^{2},\lambda(K_{Y/T}(F^{*}\xi))}^{2} = O\left( \log\log(| t |^{-2}) \right).
\end{equation}
Similarly, since the degenerate K\"ahler form $\kappa_{Y}$ on $Y$ is defined as $\kappa_{Y} = F^{*}\kappa_{X} + f^{*}\kappa_{T}$, 
the $L^{2}$-metric on $\lambda(K_{Y/T}(F^{*}\xi))$ is given by $\| \cdot \|'_{L^{2},\lambda(K_{Y/T}(F^{*}\xi))}$.
By Theorem~\ref{thm:str:sing:L2} (3), as $t \to 0$, we have
\begin{equation}
\label{eqn:estimate:L2:metric:2}
\log\|\sigma(t)\|^{\prime 2}_{L^{2},\lambda(K_{Y/T}(F^{*}\xi))} = O\left( \log\log(| t |^{-2}) \right).
\end{equation}
\par
By \eqref{eqn:comparison:Quillen:1}, \eqref{eqn:asymptotics:torsion:1}, \eqref{eqn:estimate:L2:metric:1}, \eqref{eqn:estimate:L2:metric:2}
and Proposition~\ref{prop:comparison:quillen}, we get
\begin{equation}
\label{eqn:sing:torsion:2}
\begin{aligned}
\,&
\log\tau(X_{\mu(t)}, K_{X_{\mu(t)}}(\xi))
\\
&=
\log\tau(Y_{t}, K_{Y_{t}}(F^{*}\xi))
-
\log\left(\frac{\|\cdot\|_{Q,\lambda(K_{Y_{t}}(F^{*}\xi))}}{\|\cdot\|'_{Q,\lambda(K_{Y_{t}}(F^{*}\xi))}}\right)^{2}
+
\log\left(\frac{\|\cdot\|_{L^{2},\lambda(K_{Y_{t}}(F^{*}\xi))}}{\|\cdot\|'_{L^{2},\lambda(K_{Y_{t}}(F^{*}\xi))}}\right)^{2}
\\
&=
( \alpha_{f, K_{Y/T}(F^{*}\xi)} + \beta_{Y/X, \xi} ) \log | t |^{2} + O\left( \log\log(| t |^{-2}) \right).
\end{aligned}
\end{equation}
Since $t=s^{\frac{1}{\deg\mu}}$, \eqref{eqn:kappa} follows from \eqref{eqn:asymp:eq:tors} and \eqref{eqn:sing:torsion:2}. 
This completes the proof.
\end{pf}

\end{document}
